\documentclass[a4paper,landscape]{article}
\usepackage[margin=1cm]{geometry}
\usepackage{graphicx}
\usepackage{array}
\usepackage{multirow}
\usepackage{adjustbox}
\usepackage{titlesec}
\usepackage{longtable}

\titleformat{\section}{\normalfont\Large\bfseries}{}{0pt}{}

\begin{document}
\vspace{10cm}\textit{   }\\\\
\begin{center}
\section*{Wöchentlicher Powerlifting-Trainingsplan}
\begin{adjustbox}{max width=\textwidth}
\begin{tabular}{|p{5cm}|p{5cm}|p{5cm}|p{5cm}|p{5cm}|}
\hline
\textbf{Montag} & \textbf{Dienstag} & \textbf{Mittwoch} & \textbf{Donnerstag} & \textbf{Freitag} \\
\hline
\textbf{Kniebeuge} & \textbf{Bankdrücken} & \textbf{Kreuzheben} & \textbf{Intensive Mobilitätstraining} & \textbf{Kniebeuge}\qquad (75-80-85) \\
1. Aufwärmsätze & 1. Aufwärmsätze & 1. Aufwärmsätze & 2 Min Bizeps und Trizeps & 1. Aufwärmsätze \\
2. 6×120 kg & 2. 6×75 kg & 2. 6×150 kg & 3× (fast alles) & 2. 5×130 kg \\
3. 4×136 kg & 3. 4×80 kg & 3. 4×160 kg & 3× (9×16 kg) & 3. 3×135 kg \\
4. 2×140 kg & 4. 2×85 kg & 4. 2×176 kg & & 4. 1×145 kg \\
\textbf{Tempo-Kniebeuge 250 und Beinbeuger liegend} & \textbf{Tempo-Bankdrücken 250 \& SZ-Stange Trizepsdrücken} & \textbf{Kreuzheben mit Pause (3 Sek) und Hammercurls} & \textbf{Trizepsdrücken mit Seil und Hammercurls} & \textbf{Bankdrücken} \qquad (75-80-85)  1. Aufwärmsätze \\
3×(3×115 kg) & 3×(3×70 kg) & 3×(3×140 kg) & 3×(8×14 kg) & 2. 5×80 kg \\
3×(9–12×60 kg) & 3×(10–12×Max) & 3×(10×16 kg) & \textbf{Einarmiges Trizepsdrücken} & 3. 3×85 kg \\
\textbf{Hip Thrusts und Beinstrecker} & \textbf{Kurzhantel-Schulterdrücken und Dips mit Zusatzgewicht} & \textbf{Klimmzüge und Langhantel-Biceps Curl} & \textbf{am Kabel, Reverse-Curls und Seitheben} & 4. 1×90 kg \\
2×(10×100 kg) & 3×(6–8×20 kg) & 2×bis Muskelversagen & 3×bis Muskelversagen & \textbf{Kreuzheben}\qquad (75-80-85) \\
2×(10×Max) & 3×(6–8×26 kg) & 2×(11×14 kg) & 3×bis Muskelversagen & 1. Aufwärmsätze \\
\textbf{Good Mornings und Beinstrecker} & \textbf{Seitheben und Drachenflaggen} & \textbf{Einarmiges Latziehen und Langhantel Unterarmcurl} & 3×(10×12 kg) & 2. 5×150 kg \\
2×(10×70 kg) & 4×(10×10 kg) & 2×(fast alles) & & 3. 3×160 kg \\
2×Dropset & & 2×(10–12 kg) & &  4. 1×170 kg  \\
 & & \textbf{Latziehen zur Brust und Unterarmcurl mit Langhantel} & & \textbf{Reverse Flys (Schulter)} 3×(bis Muskelversagen×10 kg) \\
 & & 2×Dropset & &   \\
 & & 2×9–13 & & \\
\hline
\end{tabular}
\end{adjustbox}
\\[1em]
\textbf{Erwartete Leistungen am Ende des Blocks\\ K: 170 kg \quad B: 105 kg \quad H: 205 kg}
\end{center}
\end{document}
